\section{Root systems}

Let $V$ be a vector space over $\R$ and $(\,,\,)$ an inner product (i.e.\ a
positive-definite symmetric bilinear form). For $\alpha \in V$,
$\alpha \neq 0$, write $\alpha\cv = \frac{2\alpha}{(\alpha, \alpha)}$, so that
$(\alpha, \alpha\cv) = 2$. Define $s_\alpha\colon V \to V$ by
$s_\alpha(v) = v - (v, \alpha\cv)\alpha$.

\begin{lemma}
$s_\alpha$ is the reflection in the hyperplane orthogonal to $\alpha$. In
particular, all but one of the eigenvalues of $s_\alpha$ are $1$, and the
remaining one is $-1$ (with eigenvector $\alpha$). Thus, $s_\alpha^2 = 1$, or
$(s_\alpha + 1)(s_\alpha - 1) = 0$, and $s_\alpha \in O(V, (\,,\,))$, the
orthogonal group of $V$ defined by the inner product $(\,,\,)$.
\end{lemma}

\begin{proof}
Write $V = \R\alpha \oplus \alpha^\perp$. Then
$s_\alpha(\alpha) = \alpha - (\alpha, \alpha\cv)\alpha = -\alpha$, and
$s_\alpha$ fixes all $v \in \alpha^\perp$.
\end{proof}

\begin{definition}
A \emph{root system} $R$ in $V$ is a finite set $R \subseteq V$ such that
\begin{enumerate}
  \item $0 \notin R$, $\R R = V$;
  \item for all $\alpha, \beta \in R$, the inner product $(\alpha, \beta\cv) \in \Z$;
  \item for all $\alpha \in R$, $s_\alpha(R) \subseteq R$. (In particular,
    $s_\alpha(\alpha) = -\alpha \in R$.)
\end{enumerate}
Moreover, $R$ is called \emph{reduced} if
\begin{enumerate}[resume]
  \item $\alpha, k\alpha \in R \implies k = \pm 1$.
\end{enumerate}
\end{definition}

\begin{example}
Let $\g$ be a semisimple Lie algebra, $\g = \at \oplus \bigoplus_{\alpha \in R} \g_\alpha$.
Then $R$ is a reduced root system in $\R R \subseteq \at^*$.
\end{example}

\begin{definition}
Let $W \subseteq GL(V)$ be the group generated by the reflections $s_\alpha$,
$\alpha \in R$; $W$ is called the \emph{Weyl group} of $R$.
\end{definition}

Observe that $|W| < \infty$: $W$ acts on $R$ by permutations, and the action
is faithful, since $R$ spans $V$ and so an element of $W$ that fixes every root
is the identity. Thus, $W$ injects into $\Sym(R)$, so
$|R| < \infty \implies |W| < \infty$.

\begin{definition}
The \emph{rank} of $R$ is the dimension of $V$.
\end{definition}

\begin{definition}
An \emph{isomorphism of root systems} $(R, V) \to (R', V')$ is a linear
bijection $\phi\colon V \to V'$ such that $\phi(R) = R'$. Note that $\phi$ is
\emph{not} required to be an isometry.
\end{definition}

\begin{definition}
If $(R, V)$ and $(R', V')$ are root systems, so is their direct sum
$(R \sqcup R', V \oplus V')$. A root system which is not isomorphic to a direct
sum of two root systems is called \emph{irreducible}.
\end{definition}

\begin{example}\leavevmode
\begin{enumerate}
  \item Rank $1$: $V = \R$, $(x, y) = xy$. $R = \{\alpha, -\alpha\}$ with
    $\alpha \neq 0$. $W = \Z/2$.
    \[
    \begin{tikzpicture}[baseline=-3pt]
      \draw[root] (0,0) -- (1.2,0) node[right] {$\alpha$};
      \draw[root] (0,0) -- (-1.2,0) node[left] {$-\alpha$};
      \node[dot,inner sep=1.3pt] at (0,0) {};
    \end{tikzpicture}
    \]

    \begin{exercise}
    This is the only rank $1$ reduced root system.
    \end{exercise}
  \item Rank $2$: $V = \R^2$ with the usual inner product.
    \begin{alist}
      \item $A_1 \times A_1$ (or $A_1 + A_1$):
        \[
        \begin{tikzpicture}[baseline=-3pt]
          \draw[root] (0,0) -- (1.2,0) node[right] {$e_1$};
          \draw[root] (0,0) -- (-1.2,0);
          \draw[root] (0,0) -- (0,1.2) node[above] {$e_2$};
          \draw[root] (0,0) -- (0,-1.2);
        \end{tikzpicture}
        \]
        This is not irreducible. The Weyl group is $\Z/2 \times \Z/2$.
      \item $A_2$: $\alpha = \alpha\cv$, $\beta = \beta\cv$,
        $(\alpha, \beta) = -1$.
        \[
        \begin{tikzpicture}[baseline=-3pt]
          \foreach \a in {0,60,...,300} \draw[root] (0,0) -- (\a:1.3);
          \node[right] at (0:1.3) {$\alpha$};
          \node[above left] at (120:1.3) {$\beta$};
        \end{tikzpicture}
        \]
        $W = S_3$ is the dihedral group of order $6$. This is the root system
        of $\sl_3$, as you should be able to establish (and we will shortly
        see).
      \item $B_2$: $\alpha = e_1$, $\beta = e_2 - e_1$; $(\alpha, \alpha) = 1$
        and $(\beta, \beta) = 2$. $\alpha$ and $\alpha + \beta$ are short
        roots; $\beta$ and $2\alpha + \beta$ are long roots.
        \[
        \begin{tikzpicture}[baseline=-3pt,scale=1.1]
          \foreach \a in {0,90,180,270} \draw[root] (0,0) -- (\a:1);
          \foreach \a in {45,135,225,315} \draw[root] (0,0) -- (\a:1.414);
          \node[right] at (0:1) {$\alpha$};
          \node[above left] at (135:1.414) {$\beta$};
          \node[above] at (90:1) {$\beta + \alpha$};
          \node[above right] at (45:1.414) {$\beta + 2\alpha$};
        \end{tikzpicture}
        \]
        $W$ is the dihedral group of order $8$. This is the root system of
        $\spl_4$ and $\so_5$.
      \item $G_2$: at this point, we clearly should look at all the dihedral
        groups, right? Unfortunately, no: the condition
        $(\alpha, \beta\cv) \in \Z$ rules out all but the dihedral group of
        order $12$. This is the root system of a Lie group, but not one of the
        classical ones (it is the automorphism group of the octonions).
        \[
        \begin{tikzpicture}[baseline=-3pt]
          \foreach \a in {0,60,...,300} \draw[root] (0,0) -- (\a:0.9);
          \foreach \a in {30,90,...,330} \draw[root] (0,0) -- (\a:1.56);
          \node[right] at (0:0.9) {$\alpha$};
          \node[left] at (150:1.56) {$\beta$};
        \end{tikzpicture}
        \]
        Here $\alpha$ is short, $\beta$ is long, $(\beta, \beta) = 3(\alpha, \alpha)$,
        and the angle between them is $150^\circ$.
    \end{alist}
\end{enumerate}
\end{example}

\begin{exercise}
All of these are indeed root systems. They are the only rank $2$ reduced root
systems, and $A_2$, $B_2$, and $G_2$ are irreducible.
\end{exercise}

\begin{lemma}
If $R$ is a root system, so is $R\cv = \setc{\alpha\cv}{\alpha \in R}$.
\end{lemma}

\begin{exercise}
Prove it.
\end{exercise}

\begin{definition}
$R$ is \emph{simply laced} if all the roots have the same length.
\end{definition}

\begin{exercise}
If $R$ is simply laced, then $(R, V) \cong (R', V')$ where $(\alpha, \alpha) = 2$
for all $\alpha \in R'$ (i.e.\ $\alpha = \alpha\cv$). [\emph{Hint:} If
irreducible, rescale; if not, break up into a direct sum of irreducibles.]
\end{exercise}

We would like to look at roots as coming from a lattice. We will see why
length $2$ is so useful.

\begin{definition}
A \emph{lattice} $L$ is a finitely generated free abelian group
($\cong \Z^l$ for some $l$) with a bilinear form $L \times L \to \Z$, such that
$(L \otimes_\Z \R, (\,,\,))$ is an inner product space. A \emph{root} of $L$ is
an $\alpha \in L$ such that $(\alpha, \alpha) = 2$. We write
$R_L = \setc{l \in L}{(l, l) = 2}$ for the roots of $L$.
\end{definition}

Note that $\alpha \in R_L \implies s_\alpha(L) \subseteq L$.

\begin{lemma}
$R_L$ is a simply laced root system in $\R R_L$.
\end{lemma}

\begin{proof}
The only non-obvious claim is $|R_L| < \infty$. However, $R_L$ is the
intersection of a compact set (the sphere of radius $\sqrt 2$) with a discrete
set (the lattice), so is finite.
\end{proof}

\begin{definition}
We say that $L$ is \emph{generated by roots} if $\Z R_L = L$. Note that then
$L$ is an \emph{even lattice}, i.e.\ $(l, l) \in 2\Z$ for all $l \in L$.
\end{definition}

\begin{example}
$L = \Z\alpha$ with $(\alpha, \alpha) = \lambda$. If $\lambda = 2$ then
$R_L = \{\pm\alpha\}$ and $L = \Z R_L$; if $\frac{k^2\lambda}{2} \neq 1$ for
all $k$ then $R_L = \emptyset$.
\end{example}

We now characterise the simply laced root systems.
\begin{enumerate}
  \item $A_n$. Consider $\Z^{n+1} = \bigoplus_{i=1}^{n+1} \Z e_i$ with
    $(e_i, e_j) = \delta_{ij}$ (the square lattice). Define
    \[
      L = \setc{l \in \Z^{n+1}}{(l, e_1 + \dots + e_{n+1}) = 0}
        = \Bigl\{\sum a_ie_i \Bigm| \sum a_i = 0\Bigr\} \cong \Z^n.
    \]
    Then $R_L = \setc{e_i - e_j}{i \neq j}$, $\#R_L = n(n + 1)$, and
    $\Z R_L = L$.

    If $\alpha = e_i - e_j$ then $s_\alpha$ applied to a vector with
    coordinates $x_1, \dots, x_{n+1}$ (in the basis $e_1, \dots, e_{n+1}$)
    swaps the $i$th and $j$th coordinates; therefore,
    $W = \langle s_{e_i - e_j}\rangle = S_{n+1}$.

    $(R_L, \R L)$ is the root system of type $A_n$. Note that $n$ is the rank
    of the root system, not the index of the associated Lie algebra.

    \begin{exercise}\leavevmode
    \begin{alist}
      \item ``You should not believe a word I say'': check these statements!
      \item Draw $L \subseteq \Z^{n+1}$ and $R_L$ for $n = 1, 2$. Check that
        $A_1$ and $A_2$ agree with the previous pictures.
      \item Show that the root system of $\sl_{n+1}$ has type $A_n$.
    \end{alist}
    \end{exercise}
  \item $D_n$. Consider the square lattice $\Z^n$. The roots are
    $R_{\Z^n} = \setc{\pm e_i \pm e_j}{i \neq j}$. Set
    \[
      L = \Z R_L = \Bigl\{ l = \sum a_ie_i \Bigm| a_i \in \Z,\ \sum a_i \text{ even} \Bigr\}.
    \]
    $s_{e_i - e_j}$ swaps the $i$th and $j$th coordinates as before;
    $s_{e_i + e_j}$ swaps them and flips the signs of both.

    $(R_L, \R R_L)$ is called $D_n$; $\#R_L = 2n(n - 1)$. The Weyl group is
    \[
      W = (\Z/2)^{n-1} \rtimes S_n,
    \]
    where $(\Z/2)^{n-1}$ is the normal subgroup of sign changes of an even
    number of coordinates.

    \begin{exercise}\leavevmode
    \begin{alist}
      \item Check these statements!
      \item $D_n$ is irreducible if $n \geq 3$.
      \item $R_{D_3} = R_{A_3}$, $R_{D_2} = R_{A_1} \times R_{A_1}$.
      \item The root system of $\so_{2n}$ has type $D_n$.
    \end{alist}
    \end{exercise}
\end{enumerate}

\lecture{12} We continue our classification of the simply laced root systems:
\begin{enumerate}
  \item $E_8$. Let
    \[
      \Gamma_n = \Bigl\{(k_1, \dots, k_n) \Bigm| \text{either all } k_i \in \Z
        \text{ or all } k_i \in \Z + \tfrac12, \text{ and } \sum k_i \in 2\Z\Bigr\}.
    \]
    Consider $\alpha = (\frac12, \frac12, \dots, \frac12) \in \Gamma_n$, and
    note that $(\alpha, \alpha) = n/4$. Thus, if $\Gamma_n$ is an even lattice,
    we must have $8 \mid n$.

    \begin{exercise}\leavevmode
    \begin{alist}
      \item $\Gamma_{8n}$ is a lattice.
      \item For $n > 1$ the roots of $\Gamma_{8n}$ form a root system of type
        $D_{8n}$.
      \item The roots in $\Gamma_8$ are
        $\setc{\pm e_i \pm e_j}{i \neq j} \cup
        \setc{\frac12(\pm e_1 \pm e_2 \pm \dots \pm e_8)}{\text{an even number of } - \text{ signs}}$.
    \end{alist}
    \end{exercise}

    The root system of $\Gamma_8$ is called the root system of type $E_8$.
    The number of roots is
    \[
      \#R_{\Gamma_8} = \binom{8}{2} \cdot 4 + 128 = 240.
    \]
    To check that you're awake: the dimension of the associated Lie algebra is
    $240 + \rank(R_{\Gamma_8}) = 248$.

    \begin{remark}
    $E_8$ is weird. Its smallest nontrivial representation is the adjoint
    representation, of dimension $248$. This is quite unlike $\so_n$, $\sl_n$,
    etc., which had dimension $O(n^2)$ and an $n$-dimensional representation.
    There isn't a good understanding of why $E_8$ is like that.
    \end{remark}

    \begin{exercise}
    Can you compute $\#W$, the order of the Weyl group of $E_8$?
    [\emph{Answer:} $2^{14} \cdot 3^5 \cdot 5^2 \cdot 7$.]
    \end{exercise}

    \begin{exercise}
    If $R$ is a root system and $\alpha \in R$, then $\alpha^\perp \cap R$ is a
    root system (in the space it spans).
    \end{exercise}
  \item $E_7$ and $E_6$. We will now apply this to $\Gamma_8$. Let
    $\alpha = \frac12(1, 1, \dots, 1)$ and $\beta = e_7 + e_8$.

    \begin{definition}
    $\alpha^\perp \cap R_{\Gamma_8}$ is a root system of type $E_7$.
    $\langle \alpha, \beta\rangle^\perp \cap R_{\Gamma_8}$ is a root system of
    type $E_6$.
    \end{definition}

    \begin{exercise}
    Show $\#R_{E_7} = 126$ and $\#R_{E_6} = 72$. Describe the corresponding
    root lattices.
    \end{exercise}

    \begin{remark}
    These lattices are reasonably natural objects. Indeed, if we took the
    associated algebraic group and its maximal torus, there is a natural
    lattice associated with it coming from the fundamental group.
    Equivalently, we can think of the group homomorphisms to (or from)
    $\C^\times$. This isn't quite the root lattice, but it's closely related
    (and we might say more about it later).
    \end{remark}
\end{enumerate}

\begin{theorem}\leavevmode
\begin{enumerate}
  \item (``ADE classification'') The complete list of irreducible simply laced
    root systems is
    \[
      A_n,\ n \geq 1, \qquad D_n,\ n \geq 4, \qquad E_6, \qquad E_7, \qquad E_8.
    \]
    No two root systems in this list are isomorphic.
  \item The remaining irreducible reduced root systems are denoted by
    \[
      B_2 = C_2, \qquad B_n,\ C_n,\ n \geq 3, \qquad F_4, \qquad G_2,
    \]
    where
    \begin{align*}
      R_{B_n} &= \setc{\pm e_i,\ \pm e_i \pm e_j}{i \neq j} \subseteq \Z^n,\\
      R_{C_n} &= \setc{\pm 2e_i,\ \pm e_i \pm e_j}{i \neq j} \subseteq \Z^n,
    \end{align*}
    with $R_{C_n}\cv = R_{B_n}$. The root system of type $B_n$ corresponds to
    $\so_{2n+1}$, and the root system of type $C_n$ corresponds to
    $\spl_{2n}$. The Weyl groups are
    \[
      W_{B_n} = W_{C_n} = (\Z/2)^n \rtimes S_n.
    \]
    The root system $F_4$ is defined as follows: let
    $Q = \setc{(k_1, \dots, k_4)}{\text{all } k_i \in \Z \text{ or all } k_i \in \Z + \frac12}$
    and take
    \[
      R_{F_4} = \setc{\alpha \in Q}{(\alpha, \alpha) = 2 \text{ or } (\alpha, \alpha) = 1}
        = \Bigl\{\pm e_i,\ \pm e_i \pm e_j\ (i \neq j),\
          \tfrac12(\pm e_1 \pm \dots \pm e_4)\Bigr\}.
    \]
\end{enumerate}
\end{theorem}

\begin{remark}
The duality between $R_{B_n}$ and $R_{C_n}$ suggests some form of duality
between $\so_{2n+1}$ and $\spl_{2n}$. We will see it later.
\end{remark}

To prove this theorem, we will first choose a good basis for $V$ (the space
spanned by the root lattice). Choose $f\colon V \to \R$ linear such that
$f(\alpha) \neq 0$ for all $\alpha \in R$.

\begin{definition}
A root $\alpha \in R$ is \emph{positive}, $\alpha \in R^+$, if $f(\alpha) > 0$.
A root $\alpha \in R$ is \emph{negative}, $\alpha \in R^-$, if $f(\alpha) < 0$.
(Equivalently, $-\alpha \in R^+$.) A root $\alpha \in R^+$ is \emph{simple},
$\alpha \in \Pi$, if $\alpha$ is not the sum of two positive roots.
\end{definition}

\begin{example}\leavevmode
\begin{enumerate}
  \item $A_n$: the roots are $R = \setc{e_i - e_j}{i \neq j}$. Choose
    \[
      f(e_1) = n + 1,\ f(e_2) = n,\ \dots,\ f(e_{n+1}) = 1.
    \]
    Then $R^+ = \setc{e_i - e_j}{i < j}$. Since $f(R) \subseteq \Z$, if
    $f(\alpha) = 1$ then $\alpha$ must be simple. Therefore,
    $\Pi \supseteq \{e_1 - e_2, \dots, e_n - e_{n+1}\}$. On the other hand, it
    is clear that the positive integer span of these is $R^+$, so in fact we
    have equality.
  \item $B_n$: $R = \setc{\pm e_i,\ \pm e_i \pm e_j}{i \neq j}$. Take
    $f(e_1) = n, \dots, f(e_n) = 1$. Then
    $R^+ = \setc{e_i,\ e_i \pm e_j}{i < j}$ and
    $\Pi = \{e_1 - e_2, \dots, e_{n-1} - e_n, e_n\}$.
  \item $C_n$: $R = \setc{\pm 2e_i,\ \pm e_i \pm e_j}{i \neq j}$.
    $R^+ = \setc{2e_i,\ e_i \pm e_j}{i < j}$ and
    $\Pi = \{e_1 - e_2, \dots, e_{n-1} - e_n, 2e_n\}$.
  \item $D_n$: $R = \setc{\pm e_i \pm e_j}{i \neq j}$,
    $\Pi = \{e_1 - e_2, \dots, e_{n-1} - e_n, e_{n-1} + e_n\}$.
  \item $E_8$: set $f(e_1) = 28 = 1 + 2 + \dots + 7$, and $f(e_i) = 9 - i$ for
    $i = 2, \dots, 8$. Then
    \[
      R^+ = \Bigl\{e_i \pm e_j\ (i < j),\ \tfrac12(e_1 \pm e_2 \pm \dots \pm e_8)
        \text{ with an even number of } - \text{ signs}\Bigr\}
    \]
    and
    \[
      \Pi = \Bigl\{\underbrace{e_2 - e_3, \dots, e_7 - e_8}_{f = 1},\
        \underbrace{\tfrac12(e_1 + e_8 - e_2 - \dots - e_7)}_{f = 1},\
        \underbrace{e_7 + e_8}_{f = 3}\Bigr\}.
    \]
\end{enumerate}
\end{example}

\begin{exercise}
Check all of these (by finding a suitable function $f$ where one is not
listed). Also find the simple roots for $E_6$, $E_7$, $F_4$, and $G_2$.
\end{exercise}

\begin{remark}
So far we have made two choices. We have chosen a maximal torus, and we have
chosen a function $f$. Neither of these matters. We may or may not prove this
about the maximal torus later; as for $f$, note that the space of functionals
$f$ is partitioned into cells within which we get the same notions of positive
roots. The Weyl group acts simply transitively on these cells, which should
show that the choice of $f$ does not affect the structure of the root system we
are deriving.
\end{remark}

\begin{proposition}\leavevmode
\begin{enumerate}
  \item $\alpha, \beta \in \Pi \implies \alpha - \beta \notin R$.
  \item $\alpha, \beta \in \Pi$, $\alpha \neq \beta \implies (\alpha, \beta\cv) \leq 0$
    (i.e.\ the angle between $\alpha$ and $\beta$ is at least $90^\circ$).
  \item Every $\alpha \in R^+$ can be written as $\alpha = \sum k_i\alpha_i$
    with $\alpha_i \in \Pi$ and $k_i \in \Z_{\geq 0}$.
  \item Simple roots are linearly independent.
  \item $\alpha \in R^+$, $\alpha \notin \Pi \implies$ there exists
    $\alpha_i \in \Pi$ such that $\alpha - \alpha_i \in R^+$.
  \item $R$ is irreducible iff $\Pi$ is indecomposable (i.e.\
    $\Pi \neq \Pi_1 \sqcup \Pi_2$ with $\Pi_1, \Pi_2$ nonempty and
    $(\Pi_1, \Pi_2) = 0$).
\end{enumerate}
\end{proposition}

\begin{exercise}
Prove this! (Either by a case-by-case analysis, which you should do in a few
cases anyway, or by finding a uniform proof from the axioms of a root system.)
\end{exercise}

We will now find another way to represent the root system. Let
$\Pi = \{\alpha_1, \dots, \alpha_l\}$, and let $a_{ij} = (\alpha_i\cv, \alpha_j)$.

\begin{definition}
$A = (a_{ij})$ is the \emph{Cartan matrix}.
\end{definition}

\begin{properties}\leavevmode
\begin{enumerate}
  \item $a_{ij} \in \Z$, $a_{ii} = 2$, $a_{ij} \leq 0$ for $i \neq j$.
  \item $a_{ij} = 0 \iff a_{ji} = 0$.
  \item $\det A > 0$.
  \item All principal minors of $A$ are positive.
\end{enumerate}
\end{properties}

\begin{proof}
The first two should be obvious from the above. To prove the third, note that
\[
  A = \begin{pmatrix}
        \frac{2}{(\alpha_1, \alpha_1)} & & 0\\
        & \ddots & \\
        0 & & \frac{2}{(\alpha_l, \alpha_l)}
      \end{pmatrix}
      \cdot \bigl((\alpha_i, \alpha_j)\bigr).
\]
The first matrix is diagonal with positive entries; the second one is the Gram
matrix of a positive definite bilinear form, and so has positive determinant.

The last property follows because any principal submatrix of $A$ (i.e.\ one
obtained by removing the same set of rows and columns) has the same form as
$A$.
\end{proof}

We will represent $A$ by a \emph{Dynkin diagram}. The vertices correspond to
simple roots. Between vertices $\alpha_i$ and $\alpha_j$ there are
$a_{ij}a_{ji}$ lines (note that this value can be $0$ or $1$ in the simply
laced diagrams, may be $2$ in $B_n$ or $C_n$, and may be $3$ in $G_2$). If
there are $2$ or $3$ lines, we put an arrow pointing towards the short root.
See figure~\ref{fig:dynkin}.

\begin{exercise}
Show the Dynkin diagrams are as we claim.
\end{exercise}

\begin{figure}[!htbp]
\centering
\begin{tikzpicture}[x=1.15cm,y=1cm,every node/.append style={font=\small}]
  % A_n
  \node[left] at (-0.6,0) {$A_n$};
  \node[node] (a1) at (0,0) {}; \node[node] (a2) at (1,0) {};
  \node[node] (a3) at (2,0) {}; \node[node] (a4) at (4,0) {};
  \draw (a1)--(a2)--(a3); \draw[dotted] (a3)--(a4);
  \node[above=2pt] at (a1) {$e_1 - e_2$}; \node[below=2pt] at (a2) {$e_2 - e_3$};
  \node[above=2pt] at (a3) {$e_3 - e_4$}; \node[below=2pt] at (a4) {$e_n - e_{n+1}$};
  % D_n
  \begin{scope}[yshift=-1.9cm]
  \node[left] at (-0.6,0) {$D_n$};
  \node[node] (d1) at (0,0) {}; \node[node] (d2) at (1,0) {};
  \node[node] (d3) at (2,0) {}; \node[node] (d4) at (3.6,0) {};
  \node[node] (d5) at (4.4,0.5) {}; \node[node] (d6) at (4.4,-0.5) {};
  \draw (d1)--(d2)--(d3); \draw[dotted] (d3)--(d4); \draw (d5)--(d4)--(d6);
  \node[above=2pt] at (d1) {$e_1 - e_2$}; \node[below=2pt] at (d2) {$e_2 - e_3$};
  \node[above=2pt] at (d3) {$e_3 - e_4$};
  \node[right=2pt] at (d5) {$e_{n-1} + e_n$}; \node[right=2pt] at (d6) {$e_{n-1} - e_n$};
  \end{scope}
  % E_8
  \begin{scope}[yshift=-3.6cm]
  \node[left] at (-0.6,0) {$E_8$};
  \foreach \i in {0,...,6} \node[node] (e\i) at (\i,0) {};
  \draw (e0)--(e6); \node[node] (eb) at (4,-0.8) {}; \draw (e4)--(eb);
  \foreach \i in {0,...,6} \node[node] at (\i,0) {};
  \end{scope}
  % E_7
  \begin{scope}[yshift=-5.1cm]
  \node[left] at (-0.6,0) {$E_7$};
  \foreach \i in {0,...,5} \node[node] (e\i) at (\i,0) {};
  \draw (e0)--(e5); \node[node] (eb) at (3,-0.8) {}; \draw (e3)--(eb);
  \foreach \i in {0,...,5} \node[node] at (\i,0) {};
  \end{scope}
  % E_6
  \begin{scope}[yshift=-6.6cm]
  \node[left] at (-0.6,0) {$E_6$};
  \foreach \i in {0,...,4} \node[node] (e\i) at (\i,0) {};
  \draw (e0)--(e4); \node[node] (eb) at (2,-0.8) {}; \draw (e2)--(eb);
  \foreach \i in {0,...,4} \node[node] at (\i,0) {};
  \end{scope}
  % B_n
  \begin{scope}[yshift=-8.3cm]
  \node[left] at (-0.6,0) {$B_n$};
  \node[node] (b1) at (0,0) {}; \node[node] (b2) at (1,0) {};
  \node[node] (b3) at (3,0) {}; \node[node] (b4) at (4,0) {};
  \draw (b1)--(b2); \draw[dotted] (b2)--(b3); \dbond{b3}{b4}
  \node[node] at (b3) {}; \node[node] at (b4) {};
  \node[above=3pt] at (b3) {$e_{n-1} - e_n$}; \node[right=4pt] at (b4) {$e_n$};
  \end{scope}
  % C_n
  \begin{scope}[yshift=-9.6cm]
  \node[left] at (-0.6,0) {$C_n$};
  \node[node] (c1) at (0,0) {}; \node[node] (c2) at (1,0) {};
  \node[node] (c3) at (3,0) {}; \node[node] (c4) at (4,0) {};
  \draw (c1)--(c2); \draw[dotted] (c2)--(c3); \dbond{c4}{c3}
  \node[node] at (c3) {}; \node[node] at (c4) {};
  \node[above=3pt] at (c3) {$e_{n-1} - e_n$}; \node[right=4pt] at (c4) {$2e_n$};
  \end{scope}
  % G_2
  \begin{scope}[yshift=-10.9cm]
  \node[left] at (-0.6,0) {$G_2$};
  \node[node] (g1) at (0,0) {}; \node[node] (g2) at (1,0) {};
  \tbond{g1}{g2} \node[node] at (g1) {}; \node[node] at (g2) {};
  \end{scope}
  % F_4
  \begin{scope}[yshift=-12.2cm]
  \node[left] at (-0.6,0) {$F_4$};
  \foreach \i in {0,...,3} \node[node] (f\i) at (\i,0) {};
  \draw (f0)--(f1); \dbond{f1}{f2} \draw (f2)--(f3);
  \foreach \i in {0,...,3} \node[node] at (\i,0) {};
  \end{scope}
\end{tikzpicture}
\caption{Dynkin diagrams (simply laced first). In $B_n$, $C_n$, $G_2$ and
$F_4$ the arrow points towards the short root.}
\label{fig:dynkin}
\end{figure}

\lecture{13}

\begin{exercise}
If $(R, V)$ is an irreducible root system with positive roots $R^+$, then there
is a \emph{unique} root $\theta \in R^+$ such that for all $\alpha \in \Pi$,
$\theta + \alpha \notin R^+$. This is called the \emph{highest root}.

At the moment, the easiest way to prove this is to examine all the root
systems we have. Later we'll give a uniform proof (and show the relationship
between $\theta$ and the adjoint representation).
\end{exercise}

\begin{example}
In $A_n$, the highest root is $e_1 - e_{n+1}$.
\end{example}

Define the ``extended Cartan matrix'' $\tilde A$ by setting
$\alpha_0 = -\theta$ (and then $\tilde A = (a_{ij})_{0 \leq i, j \leq l}$ with
$a_{ij} = (\alpha_i\cv, \alpha_j) = \frac{2(\alpha_i, \alpha_j)}{(\alpha_i, \alpha_i)}$).
The associated diagram is called the ``extended'' or ``affine'' Dynkin diagram.

\begin{example}
$A_1$ has $A = (2)$ and
$\tilde A = \bigl(\begin{smallmatrix} 2 & -2 \\ -2 & 2 \end{smallmatrix}\bigr)$.

$A_n$ has
\[
  A = \begin{pmatrix}
        2 & -1 & & \\
        -1 & 2 & \ddots & \\
        & \ddots & \ddots & -1\\
        & & -1 & 2
      \end{pmatrix}, \qquad
  \tilde A = \begin{pmatrix}
        2 & -1 & 0 & & -1\\
        -1 & 2 & -1 & & \\
        0 & -1 & 2 & \ddots & 0\\
        & & \ddots & \ddots & -1\\
        -1 & & 0 & -1 & 2
      \end{pmatrix}.
\]
\end{example}

Notice that $\tilde A$ satisfies almost the same conditions as the Cartan
matrix, i.e.
\begin{enumerate}
  \item $a_{ij} \in \Z$, $a_{ij} \leq 0$ for $i \neq j$, $a_{ii} = 2$;
  \item $a_{ij} = 0 \iff a_{ji} = 0$;
  \item $\det \tilde A = 0$, and every proper principal minor of $\tilde A$ is
    positive.
\end{enumerate}
The third property follows since $\tilde A$ is related to the Gram matrix for
$\alpha_0, \dots, \alpha_l$, which are not linearly independent.

\begin{exercise}
Write down $\tilde A$ and the extended Dynkin diagram for all types. To get you
started, in $B_n$ we have $\alpha_0 = -e_1 - e_2$, in $C_n$ we have
$\alpha_0 = -2e_1$, and in $D_n$ we have $\alpha_0 = -e_1 - e_2$.
\end{exercise}

\begin{exercise}
Show that $\tilde A_n^{(2)}$ (``twisted $A_n$'') also has determinant $0$. (See
figure~\ref{fig:affine}.)
\end{exercise}

\begin{exercise}
The Dynkin diagram of $A^T$ is the Dynkin diagram of $A$ with the arrows
reversed.
\end{exercise}

These diagrams (plus their transposes) are (almost) all the non-garbage Lie
algebras out there. (We'll be slightly more precise about the meaning of
``non-garbage'' later.)

\begin{theorem}
An irreducible (i.e., connected) Dynkin diagram is one of $A_n$, $B_n$, $C_n$,
$D_n$, $E_6$, $E_7$, $E_8$, $F_4$, and $G_2$.
\end{theorem}

\begin{proof}\leavevmode
\begin{enumerate}
  \item First, we classify the rank $2$ Dynkin diagrams, i.e.\ their Cartan
    matrices:
    \[
      A = \begin{pmatrix} 2 & -a \\ -b & 2 \end{pmatrix}
    \]
    with $\det A = 4 - ab > 0$. This leaves the options $(a, b) = (0, 0)$
    ($A_1 \times A_1$), $(1, 1)$ ($A_2$), $(2, 1)$ or $(1, 2)$ ($B_2$), and
    $(3, 1)$ or $(1, 3)$ ($G_2$).
  \item Observe that any subdiagram of a Dynkin diagram is a Dynkin diagram. In
    particular, since the affine Dynkin diagrams have determinant $0$, they
    are not subdiagrams of a Dynkin diagram.
  \item A Dynkin diagram contains no cycles. Indeed, let
    $\alpha_1, \dots, \alpha_n$ be distinct simple roots, and consider
    $\alpha = \sum \alpha_i/\sqrt{(\alpha_i, \alpha_i)}$. Then
    \[
      0 < (\alpha, \alpha) = n + \sum_{i<j}
        \frac{2(\alpha_i, \alpha_j)}{\sqrt{(\alpha_i, \alpha_i)(\alpha_j, \alpha_j)}}
        = n - \sum_{i<j} \sqrt{a_{ij}a_{ji}}.
    \]
    Therefore, $\sum_{i<j} \sqrt{a_{ij}a_{ji}} < n$. On the other hand, if we
    had a cycle on $n$ simple roots, it would have at least $n$ edges, giving
    $\sum_{i<j} \sqrt{a_{ij}a_{ji}} \geq n$, a contradiction.
  \item If a diagram is simply laced, it is $A$, $D$, or $E$.

    Suppose it's not of type $A$. Since $\tilde D_4$ is not a subdiagram, any
    branchings are at most $3$-way; and since $\tilde D_n$ is not a subdiagram
    for $n > 4$, there is exactly one branch point. Therefore, the diagram is
    T-shaped, with $p$, $q$, and $r$ vertices in each leg (counting the branch
    point in each). We call this shape $T_{p,q,r}$; e.g., $E_8 = T_{5,3,2}$.

    \begin{exercise}
    Finish the proof in one of two ways:
    \begin{alist}
      \item by using the fact that $\tilde E_6$, $\tilde E_7$, $\tilde E_8$ are
        not subdiagrams of a Dynkin diagram; or
      \item by showing $\det T_{p,q,r} = pq + pr + qr - pqr$, and concluding
        that $D_n$, $E_6$, $E_7$, and $E_8$ are the only possibilities. (E.g.:
        $\det E_8 = 15 + 10 + 6 - 30 = 1$.)
    \end{alist}
    \end{exercise}
  \item We now deal with the non-simply-laced diagrams.

    \begin{exercise}
    If there is a triple bond in the diagram, then the Dynkin diagram is
    $G_2$. (Note that we took care of one of the possibilities in
    $\tilde G_2$; you need to check the possibility of a double and a triple
    bond, and of two triple bonds.)
    \end{exercise}

    Finally, if there is a double bond in the diagram, then by $\tilde C_n$
    and $\tilde A_n^{(2)}$ there's only one of them, and by $\tilde B_n$ the
    diagram does not branch. If the double bond is in the middle, we have
    $F_4$ (since we don't have $\tilde F_4$); otherwise, we have $B_n$ or
    $C_n$. \qedhere
\end{enumerate}
\end{proof}

\begin{figure}[!htbp]
\centering
\begin{tikzpicture}[x=1.15cm,y=1cm,every node/.append style={font=\small}]
  % tilde A_n
  \node[left] at (-0.6,0) {$\tilde A_n$};
  \node[node] (a1) at (0,0) {}; \node[node] (a2) at (1,0) {};
  \node[node] (a3) at (2,0) {}; \node[node] (a4) at (4,0) {};
  \node[aff] (a0) at (2,1) {};
  \draw (a1)--(a2)--(a3); \draw[dotted] (a3)--(a4); \draw (a1)--(a0)--(a4);
  \foreach \v in {a1,a2,a3,a4} \node[node] at (\v) {}; \node[aff] at (a0) {};
  % tilde D_n
  \begin{scope}[yshift=-1.9cm]
  \node[left] at (-0.6,0) {$\tilde D_n$};
  \node[node] (d1) at (0,0.5) {}; \node[aff] (d0) at (0,-0.5) {};
  \node[node] (d2) at (0.8,0) {}; \node[node] (d3) at (1.8,0) {};
  \node[node] (d4) at (3.6,0) {};
  \node[node] (d5) at (4.4,0.5) {}; \node[node] (d6) at (4.4,-0.5) {};
  \draw (d1)--(d2)--(d0); \draw (d2)--(d3); \draw[dotted] (d3)--(d4);
  \draw (d5)--(d4)--(d6);
  \foreach \v in {d1,d2,d3,d4,d5,d6} \node[node] at (\v) {}; \node[aff] at (d0) {};
  \end{scope}
  % tilde E_8
  \begin{scope}[yshift=-3.6cm]
  \node[left] at (-0.6,0) {$\tilde E_8$};
  \foreach \i in {0,...,7} \node[node] (e\i) at (\i,0) {};
  \draw (e0)--(e7); \node[node] (eb) at (5,-0.8) {}; \draw (e5)--(eb);
  \foreach \i in {1,...,7} \node[node] at (\i,0) {}; \node[aff] at (e0) {};
  \end{scope}
  % tilde E_7
  \begin{scope}[yshift=-5.1cm]
  \node[left] at (-0.6,0) {$\tilde E_7$};
  \foreach \i in {0,...,6} \node[node] (e\i) at (\i,0) {};
  \draw (e0)--(e6); \node[node] (eb) at (3,-0.8) {}; \draw (e3)--(eb);
  \foreach \i in {0,...,5} \node[node] at (\i,0) {}; \node[aff] at (e6) {};
  \end{scope}
  % tilde E_6
  \begin{scope}[yshift=-6.6cm]
  \node[left] at (-0.6,0) {$\tilde E_6$};
  \foreach \i in {0,...,4} \node[node] (e\i) at (\i,0) {};
  \draw (e0)--(e4); \node[node] (eb) at (2,-0.8) {}; \node[aff] (ec) at (2,-1.6) {};
  \draw (e2)--(eb)--(ec);
  \foreach \i in {0,...,4} \node[node] at (\i,0) {}; \node[node] at (eb) {}; \node[aff] at (ec) {};
  \end{scope}
  % tilde B_n
  \begin{scope}[yshift=-9.1cm]
  \node[left] at (-0.6,0) {$\tilde B_n$};
  \node[node] (b1) at (0,0.5) {}; \node[aff] (b0) at (0,-0.5) {};
  \node[node] (b2) at (0.8,0) {}; \node[node] (b3) at (3,0) {};
  \node[node] (b4) at (4,0) {};
  \draw (b1)--(b2)--(b0); \draw[dotted] (b2)--(b3); \dbond{b3}{b4}
  \foreach \v in {b1,b2,b3,b4} \node[node] at (\v) {}; \node[aff] at (b0) {};
  \end{scope}
  % tilde C_n
  \begin{scope}[yshift=-10.5cm]
  \node[left] at (-0.6,0) {$\tilde C_n$};
  \node[aff] (c0) at (0,0) {}; \node[node] (c1) at (1,0) {};
  \node[node] (c2) at (2,0) {}; \node[node] (c3) at (4,0) {};
  \node[node] (c4) at (5,0) {};
  \dbond{c0}{c1} \draw (c1)--(c2); \draw[dotted] (c2)--(c3); \dbond{c4}{c3}
  \foreach \v in {c1,c2,c3,c4} \node[node] at (\v) {}; \node[aff] at (c0) {};
  \end{scope}
  % tilde G_2
  \begin{scope}[yshift=-11.8cm]
  \node[left] at (-0.6,0) {$\tilde G_2$};
  \node[aff] (g0) at (0,0) {}; \node[node] (g1) at (1,0) {};
  \node[node] (g2) at (2,0) {};
  \draw (g0)--(g1); \tbond{g1}{g2}
  \node[node] at (g1) {}; \node[node] at (g2) {}; \node[aff] at (g0) {};
  \end{scope}
  % tilde F_4
  \begin{scope}[yshift=-13.1cm]
  \node[left] at (-0.6,0) {$\tilde F_4$};
  \node[aff] (f0) at (0,0) {};
  \foreach \i in {1,...,4} \node[node] (f\i) at (\i,0) {};
  \draw (f0)--(f1)--(f2); \dbond{f2}{f3} \draw (f3)--(f4);
  \foreach \i in {1,...,4} \node[node] at (f\i) {}; \node[aff] at (f0) {};
  \end{scope}
  % tilde A_n^(2)
  \begin{scope}[yshift=-14.4cm]
  \node[left] at (-0.6,0) {$\tilde A_n^{(2)}$};
  \node[aff] (t0) at (0,0) {}; \node[node] (t1) at (1,0) {};
  \node[node] (t2) at (2,0) {}; \node[node] (t3) at (4,0) {};
  \node[node] (t4) at (5,0) {};
  \dbond{t0}{t1} \draw (t1)--(t2); \draw[dotted] (t2)--(t3); \dbond{t3}{t4}
  \foreach \v in {t1,t2,t3,t4} \node[node] at (\v) {}; \node[aff] at (t0) {};
  \end{scope}
\end{tikzpicture}
\caption{Affine Dynkin diagrams. The hatched vertex is the added root
$\alpha_0$.}
\label{fig:affine}
\end{figure}

\begin{exercise}
Compute the determinant of all the Cartan matrices.
\end{exercise}

\begin{example}
\[
  A_1\colon \det(2) = 2, \qquad
  A_2\colon \det \begin{pmatrix} 2 & -1 \\ -1 & 2 \end{pmatrix} = 3, \qquad
  A_3\colon \det = 4, \qquad \dots
\]
This comes out of the following: $SL_{n+1} = \setc{X}{\det X = 1}$ has centre
consisting of the $(n+1)$th roots of unity (i.e.\ a cyclic group of order
$n + 1$). Note that the order of the centre is equal to the determinant of the
Cartan matrix. This is true in general: the order of the centre of a simply
connected algebraic group with Lie algebra $\g$ and associated Cartan matrix
$A$ is equal to $\det A$.

In particular, this means that the simply connected algebraic group attached
to $E_8$ has no centre.
\end{example}
